% Reader-facing terminology insert for the future Version 2 manuscript.
% This file is intentionally not included by main.tex on this branch.

\subsection*{Conventions and terminology}

All graphs are finite, simple, and undirected. For a graph \(G\) and
\(S\subseteq V(G)\), write \(G[S]\) for the subgraph induced by \(S\). A set
\(S\) is \emph{independent} if \(G[S]\) is edgeless and is a \emph{clique} if
\(G[S]\) is complete. By a vertex partition we mean a family of nonempty,
pairwise disjoint sets whose union is \(V(G)\).

A \emph{cocoloring} of \(G\) is a vertex partition in which every class is an
independent set or a clique. The \emph{cochromatic number} \(\zeta(G)\) is the
minimum number of classes in a cocoloring. This is the standard invariant
introduced by Lesniak and Straight \citep{lesniak-straight-1977}; the equivalent
formulation says that every class induces an empty or complete graph. In the
language of generalized chromatic numbers,
\[
  \zeta(G)=\chi_{\mathcal P_{\mathrm{co}}}(G),
  \qquad
  \mathcal P_{\mathrm{co}}=\{K_m,\overline{K_m}:m\ge0\},
\]
where \(\mathcal P_{\mathrm{co}}\) is hereditary
\citep{scheinerman-1992,bollobas-thomason-1995}.

We write \(G(n,p)\) for the simple labelled random graph on
\([n]=\{1,\ldots,n\}\) in which the \(\binom n2\) edges occur independently
with probability \(p\). An event holds \emph{with high probability} if its
probability tends to one as \(n\to\infty\). Throughout the paper \(p=1/2\).

For nonnegative integers \(x,r\), the falling factorial is
\[
  (x)_r=x(x-1)\cdots(x-r+1),
  \qquad (x)_0=1,
\]
with the total convention \((x)_r=0\) when \(r>x\). This convention makes all
finite counting identities valid without a separate feasibility clause.

\paragraph{Signed witnesses.}
A \emph{signed cocoloring witness} is a vertex partition together with a mark
\(I\) or \(K\) on every class. An \(I\)-marked class is required to be
independent and a \(K\)-marked class is required to be a clique. The marks are
auxiliary counting data, not a new graph invariant. In the four-size profile
used below every class has size at least two for all sufficiently large \(n\),
so a realized class cannot satisfy both marks. Forgetting the marks therefore
recovers an ordinary cocoloring witness without multiplicity.

\paragraph{Overlap and configuration model.}
For ordered partitions \((V_a)_a\) and \((W_b)_b\), their overlap matrix is the
contingency table
\[
  r_{ab}=|V_a\cap W_b|.
\]
Equivalently, give row slot \(a\) exactly \(s_a\) labelled stubs and column slot
\(b\) exactly \(t_b\) labelled stubs, and choose a uniform perfect matching
between the two stub sets. The number of matched pairs in cell \((a,b)\) is
\(r_{ab}\). This bipartite configuration model is a stub-level matching and
may have several matched pairs between the same two slot types; it is not the
simple support graph defined next.

\paragraph{Support graph and binary cycle space.}
Let \(H(r)\) be the simple bipartite graph whose edges are the cells
\((a,b)\) satisfying \(r_{ab}\ge2\), with isolated slot vertices omitted. Let
\(c(H)\) denote the number of connected components, with \(c(\varnothing)=0\),
and put
\[
  \beta(H)=|E(H)|-|V(H)|+c(H).
\]
An \emph{even subgraph} is an edge set \(F\subseteq E(H)\) for which every
vertex has even degree in \((V(H),F)\). The even edge sets form the binary cycle
space under symmetric difference, and \(\beta(H)\) is its dimension. Hence
\[
  \#\{F\subseteq E(H):\deg_F(v)\text{ is even for every }v\}
  =2^{\beta(H)}.
\]

\paragraph{Manuscript-specific bookkeeping.}
The terms \emph{deficit coordinate}, \emph{demand table}, \emph{high cell},
\emph{physical high skeleton}, \emph{block support}, \emph{endpoint table},
\emph{local reward}, and \emph{residual attachment} are auxiliary terms defined
in the sections where they are used. They do not denote additional standard
graph invariants. In particular, a physical skeleton is a partial matching of
individual row and column stubs, whereas its block support is a matching of row
and column class types.
